%% ============================================================
%%  quickstart.tex — 极简入门示例（单文件可编译）
%% ------------------------------------------------------------
%%  最少代码量演示 resume-pro 的核心原语：
%%      - \makeheader      头部
%%      - \sectionTitle    带图标的章节标题
%%      - \educationItem   教育条目
%%      - \experienceItem  工作 / 实习条目
%%      - \projectCard     项目卡片
%%      - \skillbar / \skilltag   技能可视化
%%  编译：xelatex quickstart.tex
%% ============================================================

\documentclass[zh]{resume-pro}

%% 主题（可换：techblue / businessblack / academicgray / ...）
\settheme{techblue}

%% 个人信息
\name{张三}
\tagline{大模型算法工程师}
\keywords{Python \textbullet{} LLM \textbullet{} PyTorch}
\phone{+86 138-0000-0000}
\email{zhangsan@example.com}
\github{zhangsan}
\location{北京}

\begin{document}

\makeheader

\sectionTitle{教育背景}{\faGraduationCap}
\educationItem{清华大学}{2020.09 -- 2024.06}{计算机科学与技术}{本科 \textbullet{} GPA 3.8/4.0}

\sectionTitle{专业技能}{\faLaptopCode}
\skillbar{Python \textbullet{} PyTorch}{9}
\skillbar{大模型微调（LoRA / RLHF）}{8}
\skilltag{LLM} \skilltag{RAG} \skilltag{Agent} \skilltag{Docker}

\sectionTitle{实习经历}{\faBriefcase}
\experienceItem{腾讯 AI Lab}{2023.06 -- 2024.03}{大模型算法实习生}{%
\begin{itemize}
  \item 参与混元模型的 SFT 与 RLHF，benchmark 平均 +5pp
  \item 推理服务优化，P99 延迟从 3.2s 降至 1.8s
\end{itemize}
}

\sectionTitle{项目经历}{\faProjectDiagram}
\projectCard{智能问答系统}{2023.06 -- 2023.12}
{基于 RAG 的企业知识库问答}
{%
\begin{itemize}
  \item 准确率 92\%，响应时间 < 100ms
  \item 实现混合检索 + 重排序
\end{itemize}
}
{\faGithub\ github.com/zhangsan/qa-bot}

\end{document}
